% =====================================================================
\chapter{Dropout y recorte del gradiente}
\label{tema:4}
% =====================================================================

Este tema reúne dos técnicas que actúan sobre el entrenamiento sin cambiar la arquitectura. El dropout introduce ruido multiplicativo en las activaciones para regularizar; el recorte del gradiente limita la longitud de cada paso para que un gradiente anómalo no destruya lo aprendido. Se demuestran las propiedades que justifican ambas: el dropout invertido es insesgado y su varianza se calcula en forma cerrada; el recorte por norma global conserva la dirección del gradiente y acota el paso. Se sigue~\cite[§7.12]{gbc}, \cite[§5.6]{d2l}, \cite[§9.6.1]{bishopdl} y las notas de CS231n~\cite{cs231n} para el dropout, con el artículo de Srivastava \emph{et al.}~\cite{srivastava}; y \cite[§8.2.4 y §10.11]{gbc} y el artículo de Pascanu, Mikolov y Bengio~\cite{pascanu} para el recorte.

\section{Dropout}

\begin{definicion}[Dropout invertido]\label{def:dropout}
Sea $h\in\R^{d}$ un vector de activaciones y $p\in[0,1)$. Durante el entrenamiento, el \emph{dropout invertido} con probabilidad $p$ sustituye $h$ por
\[
\tilde h_i=\frac{m_i\,h_i}{1-p},\qquad m_i\sim\Bern(1-p)\ \text{independientes},\qquad i=1,\dots,d.
\]
En evaluación se usa $\tilde h=h$.
\end{definicion}

Cada unidad se apaga con probabilidad $p$ y las que sobreviven se amplifican por $1/(1-p)$. La variante original~\cite{srivastava} no amplificaba durante el entrenamiento y, en cambio, multiplicaba los pesos por $1-p$ en la evaluación; la variante invertida traslada esa corrección al entrenamiento y deja intacto el código de inferencia~\cite{cs231n}.

\begin{proposicion}[Esperanza y varianza]\label{prop:dropout}
Con la notación de la \cref{def:dropout}, para $h$ fijo,
\[
\E[\tilde h_i]=h_i,\qquad \Var(\tilde h_i)=\frac{p}{1-p}\,h_i^{2}.
\]
\end{proposicion}

\begin{corolario}[Dropout antes de una capa lineal]\label{cor:dropout-lineal}
Sea $y=w\T\tilde h$ con $w\in\R^{d}$ fijo. Entonces $\E[y]=w\T h$ y $\Var(y)=\frac{p}{1-p}\sum_{i=1}^{d}w_i^{2}h_i^{2}$.
\end{corolario}

El corolario muestra que, para una capa lineal, evaluar sin máscara da exactamente la salida media sobre todas las máscaras. Goodfellow \emph{et al.}~\cite[§7.12]{gbc} interpretan el dropout como el entrenamiento simultáneo de los $2^{d}$ subredes que se obtienen al eliminar unidades, con parámetros compartidos; la regla de evaluación sin máscara, llamada \emph{regla de escalado de pesos}, es exacta para modelos sin no linealidades ocultas, como el anterior, y solo aproximada en redes profundas no lineales.

\addtocounter{definicion}{1}%

\addtocounter{definicion}{1}%

\section{Recorte del gradiente por norma global}

\begin{definicion}[Recorte por norma global]\label{def:recorte}
Sean $g^{(1)},\dots,g^{(K)}$ los gradientes de los $K$ tensores de parámetros de un modelo y sea $g\in\R^{n}$ su concatenación. Dado un umbral $c>0$, el \emph{recorte por norma global} sustituye cada $g^{(k)}$ por
\[
\hat g^{(k)}=\alpha\,g^{(k)},\qquad \alpha=\min\Bigl(1,\ \frac{c}{\norm{g}_2}\Bigr),
\]
con el convenio $\alpha=1$ si $g=0$.
\end{definicion}

\begin{proposicion}[Propiedades del recorte]\label{prop:recorte}
Sea $\hat g=\alpha g$ el gradiente recortado.
\begin{enumerate}
  \item $\hat g$ tiene la dirección y el sentido de $g$, y $\norm{\hat g}_2=\min(\norm{g}_2,\,c)$.
  \item Un paso de descenso $\theta\leftarrow\theta-\eta\hat g$ cumple $\norm{\Delta\theta}_2\le\eta c$.
  \item Si $\hat g\ne 0$, $\hat g$ es una dirección de descenso: $\langle\nabla\Loss,-\hat g\rangle<0$ cuando $\hat g$ se calcula con el gradiente exacto.
  \item Recortar cada tensor por separado con el mismo umbral conserva, en general, la dirección de cada $g^{(k)}$ pero no la de $g$.
\end{enumerate}
\end{proposicion}

\addtocounter{definicion}{1}%

El recorte lo propusieron Pascanu, Mikolov y Bengio~\cite[§3.2]{pascanu} para redes recurrentes, donde el \cref{cor:pascanu} permite gradientes que explotan; en un acantilado de la superficie de pérdida~\cite[§8.2.4]{gbc}, el gradiente apunta en una dirección razonable pero tiene una longitud desproporcionada, y el recorte conserva lo primero y corrige lo segundo. En PyTorch se aplica con \texttt{torch.nn.utils.clip\_grad\_norm\_} entre \texttt{backward()} y el paso del optimizador.

\addtocounter{definicion}{1}%

